\section{Anhang: Offene Arbeiten am Bericht (Referenz)}

Die folgende Liste dient als projektinterne Checkliste für die noch ausstehenden Arbeiten an diesem Bericht.

\begin{itemize}
    \item[\checkmark] wissenschaftliche Quellen und Literaturverweise ergänzen,
    \item[\checkmark] Abbildungen und Tabellen nummerieren und im Text referenzieren,
    \item[\checkmark] Literaturverzeichnis mit allen zitierten Arbeiten vervollständigen,
    \item[\checkmark] Sensor-Setup-Abbildung einfügen,
    \item[\checkmark] qualitative Detektionsbeispiele einfügen (\texttt{os1} gegenüber \texttt{merged}, Standbild aus dem Video),
    \item[\checkmark] Baselinevergleich vortrainiert gegenüber finetunet ergänzen,
    \item[\checkmark] mAP der bewegten Objekte und experimentweises Raster ergänzen,
    \item[\checkmark] Bewegt-mAP foldweise als Mittel \(\pm\) Standardabweichung angeben,
    \item[\checkmark] alle zentralen Ergebnisse zusätzlich unter der offiziellen Variante berichten,
    \item[\checkmark] AP\textsubscript{40} und AP\textsubscript{50} tabellieren,
    \item[\checkmark] Personen-Einbruch von \texttt{os0} gegen die Foldwerte verifizieren,
    \item[\checkmark] Trainingsparameter vervollständigen,
    \item[\checkmark] Merge-Zeitverteilung einordnen.
\end{itemize}

\noindent Offen sind noch:

\begin{itemize}
    \item[\(\ast\)] \textbf{Sensorspezifikation OS0 und OS1 ergänzen} (Modell, Kanalzahl,
          vertikales Sichtfeld, Reichweite, Montagehöhe). Der Bericht beschreibt an
          mehreren Stellen Asymmetrien zwischen den Einzelsensoren, ohne eine
          physikalische Erklärung anzubieten. Höchster verbleibender Erkenntnisgewinn.
    \item Experimentweises Raster auch für AP\textsubscript{60} erstellen. Die Werte liegen
          in \texttt{CROSS\_VALIDATION\_RESULTS.md} vor und würden den Fusionsvorteil
          auch experimentweise belegen.
    \item Angaben zur Sensorsynchronisation ergänzen (Verfahren, maximaler Zeitversatz).
          Ein Zeitversatz wirkt bei bewegten Objekten unmittelbar auf AP\textsubscript{60}.
    \item ICP-Parameter und Registrierungsgüte dokumentieren (Voxelgröße,
          Korrespondenzdistanz, Abbruchkriterium, Fitness bzw.\ RMSE).
    \item Kriterium dokumentieren, ab welcher Sichtbarkeit ein Objekt in einer Ansicht
          annotiert wurde. Davon hängen die ansichtsspezifischen GT-Zahlen ab.
    \item Datum der Definition der labelbasierten Bewegt-\slash Statisch-Variante nachtragen.
    \item Datenstand der Ablationen aus Abschnitt~\ref{sec:gtsampling} und
          \ref{sec:ursachen} klären; gegebenenfalls auf dem korrigierten Stand wiederholen.
    \item Architekturschema von PointPillars einfügen,
    \item Tabelle mit Frameanzahl und Instanzen je Experiment und Klasse ergänzen,
    \item prüfen, ob eine Eigenständigkeitserklärung und eine Aufteilung der
          Beiträge beider Teammitglieder verlangt sind.
\end{itemize}

\noindent Mit \(\ast\) markierte Punkte betreffen die Belastbarkeit von Aussagen im Ergebniskapitel.
